% GENERATED FILE. Edit canonical lessons, then run npm run generate:books.
% canonical-source-hash: fnv1a64:f1c406354c1e80b7
% canonical-lessons: TA-C191-kataikkarar, TA-C191-vatikkaiyalar, TA-C191-kacalar, TA-C191-taiyalkarar, TA-C191-cillarai

\chapter{Shopkeeper, Customer, Cashier, Tailor, Small change}
\label{ch:ta-shopkeeper-customer-cashier-tailor-small-change}
\hlchaptermodality{191}

\begin{chapteropening}
\textbf{By the end of this chapter:} \emph{I can say a shopkeeper, a customer, a cashier, a tailor and small change.}

Complete the last lesson of chapter 191: I can say a shopkeeper, a customer, a cashier, a tailor and small change.
\end{chapteropening}

\section[\texorpdfstring{\ta{கடைக்காரர்} (kaṭaikkārar)}{கடைக்காரர் (kaṭaikkārar)}]{\ta{கடைக்காரர்} (kaṭaikkārar) — a shopkeeper}

\label{lesson:TA-C191-kataikkarar}



\begin{quote}
Before the new one: say the Tamil for a grinding stone, then the Tamil for a tank.
\end{quote}

\subsection*{You'll want to know: \ta{கடைக்காரர்}}
\textbf{\ta{கடைக்காரர்}} — \emph{kaṭaikkārar} — \textquotedblleft{}a shopkeeper\textquotedblright{}.

\begin{grammarlens}[title={What you've built}]
One more word for this chapter.
\end{grammarlens}

\subsection*{Your turn}
\begin{itemize}
\raggedright
  \item \emph{Say it:} \emph{kaṭaikkārar}
  \item \emph{Say it:} \emph{kaṭaikkārar}, once more
  \item \emph{Say it:} say \emph{kaṭaikkārar}
  \item \emph{Recall:} say the Tamil for a grinding stone, then the Tamil for a tank, then say \emph{kaṭaikkārar} again
\end{itemize}

\begin{tcolorbox}[breakable,colback=teal!4,colframe=teal!35!black,title={Before you move on}]
What does \ta{கடைக்காரர்} mean? (\textquotedblleft{}A shopkeeper\textquotedblright{}.) Say it once more.
\end{tcolorbox}
\section[\texorpdfstring{\ta{வாடிக்கையாளர்} (vāṭikkaiyāḷar)}{வாடிக்கையாளர் (vāṭikkaiyāḷar)}]{\ta{வாடிக்கையாளர்} (vāṭikkaiyāḷar) — a customer}

\label{lesson:TA-C191-vatikkaiyalar}



\begin{quote}
Before the new one: say the Tamil for a tank, then the Tamil for a shopkeeper.
\end{quote}

\subsection*{You'll want to know: \ta{வாடிக்கையாளர்}}
\textbf{\ta{வாடிக்கையாளர்}} — \emph{vāṭikkaiyāḷar} — \textquotedblleft{}a customer\textquotedblright{}.

\begin{grammarlens}[title={What you've built}]
One more word for this chapter.
\end{grammarlens}

\subsection*{Your turn}
\begin{itemize}
\raggedright
  \item \emph{Say it:} \emph{vāṭikkaiyāḷar}
  \item \emph{Say it:} \emph{vāṭikkaiyāḷar}, once more
  \item \emph{Say it:} say \emph{vāṭikkaiyāḷar}
  \item \emph{Recall:} say the Tamil for a tank, then the Tamil for a shopkeeper, then say \emph{vāṭikkaiyāḷar} again
\end{itemize}

\begin{tcolorbox}[breakable,colback=teal!4,colframe=teal!35!black,title={Before you move on}]
What does \ta{வாடிக்கையாளர்} mean? (\textquotedblleft{}A customer\textquotedblright{}.) Say it once more.
\end{tcolorbox}
\section[\texorpdfstring{\ta{காசாளர்} (kācāḷar)}{காசாளர் (kācāḷar)}]{\ta{காசாளர்} (kācāḷar) — a cashier}

\label{lesson:TA-C191-kacalar}



\begin{quote}
Before the new one: say the Tamil for a shopkeeper, then the Tamil for a customer.
\end{quote}

\subsection*{You'll want to know: \ta{காசாளர்}}
\textbf{\ta{காசாளர்}} — \emph{kācāḷar} — \textquotedblleft{}a cashier\textquotedblright{}.

\begin{grammarlens}[title={What you've built}]
One more word for this chapter.
\end{grammarlens}

\subsection*{Your turn}
\begin{itemize}
\raggedright
  \item \emph{Say it:} \emph{kācāḷar}
  \item \emph{Say it:} \emph{kācāḷar}, once more
  \item \emph{Say it:} say \emph{kācāḷar}
  \item \emph{Recall:} say the Tamil for a shopkeeper, then the Tamil for a customer, then say \emph{kācāḷar} again
\end{itemize}

\begin{tcolorbox}[breakable,colback=teal!4,colframe=teal!35!black,title={Before you move on}]
What does \ta{காசாளர்} mean? (\textquotedblleft{}A cashier\textquotedblright{}.) Say it once more.
\end{tcolorbox}
\section[\texorpdfstring{\ta{தையல்காரர்} (taiyalkārar)}{தையல்காரர் (taiyalkārar)}]{\ta{தையல்காரர்} (taiyalkārar) — a tailor}

\label{lesson:TA-C191-taiyalkarar}



\begin{quote}
Before the new one: say the Tamil for a customer, then the Tamil for a cashier.
\end{quote}

\subsection*{You'll want to know: \ta{தையல்காரர்}}
\textbf{\ta{தையல்காரர்}} — \emph{taiyalkārar} — \textquotedblleft{}a tailor\textquotedblright{}.

\begin{grammarlens}[title={What you've built}]
One more word for this chapter.
\end{grammarlens}

\subsection*{Your turn}
\begin{itemize}
\raggedright
  \item \emph{Say it:} \emph{taiyalkārar}
  \item \emph{Say it:} \emph{taiyalkārar}, once more
  \item \emph{Say it:} say \emph{taiyalkārar}
  \item \emph{Recall:} say the Tamil for a customer, then the Tamil for a cashier, then say \emph{taiyalkārar} again
\end{itemize}

\begin{tcolorbox}[breakable,colback=teal!4,colframe=teal!35!black,title={Before you move on}]
What does \ta{தையல்காரர்} mean? (\textquotedblleft{}A tailor\textquotedblright{}.) Say it once more.
\end{tcolorbox}
\section[\texorpdfstring{\ta{சில்லறை} (cillaṟai)}{சில்லறை (cillaṟai)}]{\ta{சில்லறை} (cillaṟai) — small change}

\label{lesson:TA-C191-cillarai}



\begin{quote}
Before the new one: say the Tamil for a cashier, then the Tamil for a tailor.
\end{quote}

\subsection*{You'll want to know: \ta{சில்லறை}}
\textbf{\ta{சில்லறை}} — \emph{cillaṟai} — \textquotedblleft{}small change\textquotedblright{}.

\begin{grammarlens}[title={What you've built}]
One more word for this chapter.
\end{grammarlens}

\subsection*{Your turn}
\begin{itemize}
\raggedright
  \item \emph{Say it:} \emph{cillaṟai}
  \item \emph{Say it:} \emph{cillaṟai}, once more
  \item \emph{Say it:} say \emph{cillaṟai}
  \item \emph{Recall:} say the Tamil for a cashier, then the Tamil for a tailor, then say \emph{cillaṟai} again
\end{itemize}

\begin{tcolorbox}[breakable,colback=teal!4,colframe=teal!35!black,title={Before you move on}]
What does \ta{சில்லறை} mean? (\textquotedblleft{}Small change\textquotedblright{}.) Say it once more.
\end{tcolorbox}
